\section{Data and Evaluation Protocol}\label{sec:data}

\subsection{EgoPoint}
We train and test on the image set of EgoPoint-Ground~\cite{egopoint}, which we call EgoPoint. Its paper reports 15,338 images:
\begin{itemize}
\item 39\% (6,035) real captures with RayNeo~V3 smart glasses (12\,MP Sony IMX681 sensor);
\item 39\% (5,976) synthetically composed;
\item 22\% (3,327) LLM-generated.
\end{itemize}
These are the paper's main-text labels; its appendix calls the last two groups ``synthetic generation'' and ``image editing''. The paper describes a random 7:2:1 train:val:test split, plus a domain-adaptive split that keeps all real captures for testing. The release we use has 15,455 images split 10,837 / 1,546 / 3,072. VistaRef~\cite{li2026vistaref} and PointVG-R~\cite{li2026pointvgr} report the same counts, so our test images are theirs, minus the exclusions below. \tbd{Ask the EgoPoint-Ground authors how this release relates to the paper's splits, and state the access terms and licence.} Each positive record shows one wearer's hand pointing at one target; negatives show no hand. Per-image source labels (real, composed or generated) are \tbd{not available in our copy; if obtained, also report results on the real subset}.

\paragraph{Conversion.} We convert each record to our label format (the three flags, hand box, fingertip and target box):
\begin{itemize}
\item Negatives become 000, and their object boxes are ignored.
\item Positives with a hand box, an annotated fingertip and a target box become 111.
\item In three test records, the hand label was attached to an object and the hand was labelled as an object. We detected these by the fingertip position and swapped the labels back.
\item We dropped 232 positives: 229 without an annotated fingertip (164 train, 22 validation, 43 test) and 3 without a hand box. All 229 are flagged as unclear in the source annotations. Spot checks showed open palms, held objects, or a hand box on the hand that was not pointing.
\end{itemize}
Table~\ref{tab:data} summarises the result. Images are stored at $224\times224$; the model squashes every frame to its square $160\times160$ input, regardless of the original aspect ratio.

\begin{table}[t]
  \caption{EgoPoint after conversion. \emph{On object}: the fingertip lies inside the target box; \emph{near}: fingertip--box distance ${\leq}\,0.08$ of the image side; \emph{far}: ${>}\,0.08$.}
  \label{tab:data}
  \footnotesize
  \begin{tabular}{lrrrrrr}
    \toprule
    & & & & \multicolumn{3}{c}{Pointing at target, by distance} \\
    \cmidrule(l){5-7}
    Split & Images & No hand (000) & Pointing at target (111) & On object & Near & Far \\
    \midrule
    Train & 10,669 & 1,535 & 9,134 & 6,511 & 1,995 & 628 \\
    Validation & 1,524 & 211 & 1,313 & 922 & 297 & 94 \\
    Test & 3,029 & 425 & 2,604 & 1,796 & 599 & 209 \\
    \midrule
    Total & 15,222 & 2,171 & 13,051 & 9,229 & 2,891 & 931 \\
    \bottomrule
  \end{tabular}
\end{table}

\paragraph{Consequences for evaluation.}
\begin{itemize}
\item \emph{Test split.} Our test split has 3,029 images, 43 positives fewer than the official one. Wherever we show numbers published on the official split, we say so; those numbers are not strictly comparable (Section~\ref{sec:main}).
\item \emph{Two scenes only.} EgoPoint contains only scenes 000 and 111. The pointing and target gates therefore never see a negative example, and we report neither their accuracy nor four-class scene accuracy.
\item \emph{Distance skew.} About 70\% of targets are touched by the fingertip, and only 7\% are far. Pooled results are dominated by touched targets, so we report every target metric by distance as well. The median distance from the fingertip to the target centre is 0.12, 0.17 and 0.32 of the image side for on-object, near and far targets (test split).
\end{itemize}

\subsection{Metrics}\label{sec:metrics}
\begin{itemize}
\item \emph{Target and hand localisation.} We use the two standard grounding metrics~\cite{egopoint,li2026vistaref,li2026pointvgr}. \emph{mIoU} is the IoU between predicted and labelled box, averaged over images, with a 95\% bootstrap confidence interval (1,000 resamples~\cite{efron1994introduction}). \emph{P@$\tau$} (Precision@$\tau$) is the share of images with IoU${\geq}\tau$, for $\tau$ = 0.3, 0.5 and 0.7. Both are given in \%. Target metrics are given overall and by distance.
\item \emph{Fingertip and direction.} Fingertip error is in pixels at the $160\times160$ input (1\,px = 0.63\% of the image side). Direction error is the angle between $\hat{\dir}$ and $\dir^\star$.
\item \emph{Hand presence.} For the hand gate, we report accuracy, missed hands, false alarms and AUROC, and the two-scene accuracy (no hand versus pointing at target), with Wilson intervals~\cite{wilson1927probable} for proportions.
\item \emph{No skipped images.} Localisation is scored on every image whose label contains the part, whatever scene the model predicted, so hard images cannot be skipped.
\end{itemize}

\paragraph{Ownership test.} To test robustness to other people's hands, we paste one or two hands into every test image. The hands come from other test images, 30\% are re-pasted copies of the wearer's own hand, and their positions are fixed by a seed. Labels stay unchanged. We report:
\begin{itemize}
\item \emph{ownership accuracy}: the share of hand images in which the predicted hand centre is nearer the wearer's hand than any pasted hand;
\item how often the predicted centre falls inside a pasted hand;
\item the share of no-hand images that stay \emph{no hand}.
\end{itemize}
These distractors follow the training distribution, so this test measures robustness to that synthetic distribution and may overstate robustness in real scenes. Section~\ref{sec:own} therefore adds real frames.

\paragraph{Model selection.} The reported checkpoint has the lowest validation loss. It was fixed before we looked at the test split, and the test split was evaluated once. \tbd{Confirm: loss weights and all hyper-parameters were set on the validation split only.}
